\documentclass[9pt,a4paper]{article}
% Scheda di correzione per l'alunno: una sola pagina A4. Nessun indicatore ANA/SVI/ARG/COM, nessun ID evidenza.
\usepackage[T1]{fontenc}
\usepackage[utf8]{inputenc}
\usepackage[italian]{babel}
\usepackage{mathpazo,microtype,amsmath}
\usepackage{siunitx}
\sisetup{output-decimal-marker={,},per-mode=symbol,group-digits=false}
\usepackage{pifont,xcolor,booktabs,array,tabularx,enumitem}
\usepackage[a4paper,margin=1.5cm]{geometry}
\setlength{\parindent}{0pt}\setlength{\parskip}{3pt}\pagestyle{empty}
\setlist[itemize]{leftmargin=1.1em,itemsep=1pt,topsep=1pt,parsep=0pt}

% Rosso e verde accesi riservati ai simboli di esito (solo in questa scheda, mai in prove e correttore).
\definecolor{okverde}{HTML}{00B050}\definecolor{kosso}{HTML}{E60000}
\newcommand{\ok}{\item[{\color{okverde}\ding{51}}]}
\newcommand{\ko}{\item[{\color{kosso}\ding{55}}]}
\newcommand{\esito}[2]{% #1 = elenco \ok, #2 = elenco \ko
  \begin{tabularx}{\linewidth}{@{}X@{\hspace{6pt}}X@{}}
    \begin{minipage}[t]{\linewidth}\textbf{\color{okverde}Corretto}\begin{itemize}[label={}] #1\end{itemize}\end{minipage} &
    \begin{minipage}[t]{\linewidth}\textbf{\color{kosso}Errori e mancanze}\begin{itemize}[label={}] #2\end{itemize}\end{minipage}
  \end{tabularx}}
\newcommand{\esercizio}[2]{\par\medskip\textbf{Esercizio #1}\hfill #2\par}

\begin{document}
{\Large\bfseries Verifica \textsc{ARGOMENTO} --- COGNOME NOME}\\
Classe CLASSE \quad Data DATA \quad Punteggio massimo 1000

\medskip
\begin{tabular}{@{}lrr@{}}\toprule
 & Punti & Massimo\\\midrule
 % una riga per esercizio: Esercizio 1 & 180 & 250 \\
 \midrule \textbf{Totale} & \textbf{TOT} & 1000\\\bottomrule
\end{tabular}
\hfill {\Large\bfseries Voto: VOTO}

% \esercizio{1}{180 / 250}
% \esito{\ok formula giusta \ok unità corrette}{\ko manca la verifica del dominio}

\medskip\textbf{Aree di miglioramento}
\begin{itemize}[label=$\rightarrow$]
 \item PUNTO 1
 \item PUNTO 2
\end{itemize}

{\small Calcolo: punti TOT, dopo la divisione per 100, arrotondati a 0,25 \(\Rightarrow\) VOTO.}
\end{document}
